\section{Introduction}\label{sec:introduction}
% P:introduction-01
Monoethanolamine (\MEA) in water is the reference solvent for \COtwo absorption, and its vapor--liquid equilibrium (VLE) has been measured extensively~\cite{Jou1995,Mamun2005,Hilliard2008,Aronu2011,Xu2011}.
Absorbed \COtwo reacts with water and amine to form protonated amine, carbamate, bicarbonate and carbonate.
The equilibrium \COtwo partial pressure at a given loading therefore depends on both the phase equilibrium and the liquid speciation.
A model can reproduce the pressure with an implausible distribution of \MEA, \MEAH, \MEACOO, \HCOthree and \COthree, so speciation measurements provide a separate test of the same thermodynamic state.

% P:introduction-02
Speciation measurements for loaded \MEA differ in what they observe.
Quantitative nuclear magnetic resonance (NMR) resolves carbamate. B\"ottinger et al.\ report the sum of molecular and protonated amine and tabulate bicarbonate, noting that carbonate is practically absent in their studied solutions~\cite[p.~134, Table~3]{Bottinger2008}; Jakobsen et al.\ use calibrated chemical shifts to separate molecular from protonated amine and bicarbonate from carbonate~\cite[pp.~9895--9896]{Jakobsen2005}.
Matin et al.\ obtained all four major species from acid and base titrations combined with a total carbon measurement~\cite{matinFacileMethodDetermination2012}.
They reported agreement with NMR for free amine, protonated amine and carbamate, and bicarbonate higher than the NMR values~\cite[pp.~6616--6617]{matinFacileMethodDetermination2012}.
Each comparison must distinguish the reported observable from the calculated observation operator, including any approximation used to relate them.

% P:introduction-03
Multi-property calibration of aqueous \MEA is established in electrolyte excess-Gibbs-energy models.
B\"ottinger et al.\ fitted carbamate formation to NMR while retaining the Bates--Pinching protonation law and previously estimated Pitzer interactions \cite[Sects.~5--6, Tables~5--6]{Bottinger2008}.
Zhang et al.\ jointly fitted electrolyte non-random two-liquid (NRTL) parameters to solubility, absorption heat, heat capacity and NMR; Akula et al.\ fitted pressure and differential heat, then compared NMR and heat capacity \cite[Sect.~3.3, Tables~9--11]{Zhang2011}\cite[Sects.~3.3 and 4.2]{Akula2023a}.
Hessen et al.\ re-estimated refined electrolyte NRTL interactions against binary properties and ternary pressures, while Aronu et al.\ used extended UNIQUAC for their 15--60 wt\% MEA measurements \cite[Sect.~4.1, Table~2]{Hessen2010}\cite[Sects.~3--4]{Aronu2011}.
Here we examine a different combination: explicit-ion ePC-SAFT with carbamate chemistry and a solvation-shell Born term, assessed against both property data and finite phase searches.

% P:introduction-07
SAFT models represent reactive absorption either through association sites or through separate chemical species.
SAFT-VR, SAFT-$\gamma$ square-well and SAFT-$\gamma$ Mie treatments derive associated species from bonded-site fractions \cite{MacDowell2010,Rodriguez2012,chremosModellingPhaseChemical2016,perdomoPredictiveGroupcontributionFramework2023}; soft-SAFT also estimates absorption enthalpy after fitting solubility \cite[Sect.~4.7]{lloretConsistentTransferableThermodynamic2017}.
Neutral-species PC-SAFT and Huang--Radosz SAFT instead combine phase fugacities with ideal reaction equilibria \cite{Nasrifar2010,Baygi2015,Najafloo2018}.
In Nasrifar and Tafazzol's water balance, hydroxide, bicarbonate and carbonate amounts contribute to the effective water amount \cite[Eq.~12]{Nasrifar2010}; these approaches do not assign separate nonideal activities to each ion.
Explicit-ion alternatives include T\'ellez-Arredondo and Medeiros' electrolyte cubic two-state association model and Neumann et al.'s combination of Helmholtz-energy equations of state with electrolyte NRTL \cite{tellez-arredondoModelingCO2H2S2013,neumannNewModelCombining2022}.
Their different species representations and fitted data are compared quantitatively in \cref{tab:literature-model-comparison}.

% P:introduction-04
Electrolyte PC-SAFT (ePC-SAFT) adds a Debye--H\"uckel contribution for charged species~\cite{Cameretti2005}, and Held et al.\ estimated ion parameters from densities and osmotic coefficients of aqueous salt solutions, with dispersion between ions of the same sign excluded~\cite{Held2014}.
With explicit ions and activity-based reaction equilibria, ePC-SAFT described acid-gas solubility in aqueous MDEA from parameters obtained from pure components and binary subsystems~\cite{Uyan2015,Wangler2018}, and was extended to competitive \ce{H2S}--\COtwo absorption~\cite{Cleeton2020}.
B\"ulow et al.\ added the Born term for the solvation Gibbs energy of ions to MDEA blends with sulfolane or an ionic liquid, which reduced the number of binary interaction parameters~\cite{Bulow2021a}.
Related work added concentration-dependent permittivity to the Born and Debye--H\"uckel terms~\cite{Bulow2020,Rueben2024} and predicted \COtwo solubility in aqueous and organic salt solutions~\cite{Schick2023}.
Figiel et al.\ modified the Born term to account for the solvation-shell volume and dielectric saturation near an ion (SSM+DS), with parameters estimated from ion solvation and salt activity data in water and alcohols at \SI{298.15}{\kelvin}~\cite{Figiel2025}; a later correction supplied a missing bracket and the ion-region permittivity in its Eq.~(7)~\cite{Figiel2026Correction}.
MDEA is a tertiary amine and forms no carbamate.
The primary amine \MEA forms carbamate, so its ion population and solvent composition change together with \COtwo loading.

% P:introduction-05
Density and phase searches test different consequences of a fitted Helmholtz energy.
Agreement with pressure and speciation need not imply accurate liquid volume, and a converged equilibrium state may still admit a lower-Gibbs-energy phase arrangement.
We compare local fits of the same nine-species model under original and restricted parameter bounds to quantify the accompanying property discrepancies.
The phase searches cover only 313.15--\SI{333.15}{\kelvin} and loading $10^{-5}$--0.5; the higher-temperature property comparisons require separate assessment of phase behavior.
This comparison addresses thermodynamic representation, with the scope and unresolved discrepancies stated in \cref{sec:results-limits}.
